\section{第6章 定积分的应用}
\subsection{定积分在几何学上的应用}
\begin{exercise}{661}{基础·曲边梯形面积}
求 $y=\sqrt x$、$y=0$ 与 $x=4$ 围成区域的面积。
\end{exercise}
\begin{exercise}{662}{基础·横向切片面积}
求由 $x=y^2$、$x=0$、$y=1$ 围成区域的面积，使用关于 $y$ 的积分。
\end{exercise}
\begin{exercise}{663}{常规·面积与旋转体积}
求 $y=x^2$ 与 $y=x$ 围成区域的面积，并求该区域绕 $x$ 轴旋转一周的体积。
\end{exercise}
\begin{exercise}{664}{基础·抛物线间面积}
求 $y=4-x^2$ 与 $y=4-2x$ 围成区域的面积，先求交点。
\end{exercise}
\begin{exercise}{665}{常规·面积反求参数}
曲线 $y=\sqrt{ax}$（$a>0$）、$x$ 轴及 $x=3$ 围成区域的面积为 $6$。求 $a$。
\end{exercise}
\begin{exercise}{666}{基础·指数曲线面积}
求 $y=e^x$、$y=0$、$x=0$、$x=1$ 围成区域的面积。
\end{exercise}
\begin{exercise}{667}{基础·有向积分与面积}
求 $y=x^2-x$ 与 $x$ 轴围成区域的面积。
\end{exercise}
\begin{exercise}{668}{常规·切线与面积}
求 $y=x^2$ 在 $(1,1)$ 处的切线，并求切线、抛物线与 $y$ 轴围成区域的面积。
\end{exercise}
\begin{exercise}{669}{基础·横向切片}
求 $x=4-y^2$ 与 $x=0$ 围成区域的面积。
\end{exercise}
\begin{exercise}{670}{基础·三角曲线面积}
求 $y=\sin x$（$0\le x\le\pi$）与 $x$ 轴围成区域的面积。
\end{exercise}
\begin{exercise}{671}{常规·对称区域面积}
求 $y=e^{-|x|}$、$y=0$、$x=-1$、$x=1$ 围成区域的面积。
\end{exercise}
\begin{exercise}{672}{基础·对数曲线面积}
求 $y=\ln x$、$y=0$ 与 $x=e$ 围成区域的面积。
\end{exercise}
\begin{exercise}{673}{常规·两三角曲线间面积}
求 $y=\sin x$、$y=\cos x$、$x=0$、$x=\frac{\pi}{2}$ 围成区域的面积。
\end{exercise}
\begin{exercise}{674}{基础·圆弓形面积}
求单位圆 $x^2+y^2\le1$ 位于 $x$ 轴上方的区域面积，用定积分表示并计算。
\end{exercise}
\begin{exercise}{675}{基础·水平直线截面积}
求 $y=4-x^2$ 与 $y=3$ 围成区域的面积。
\end{exercise}
\begin{exercise}{676}{基础·椭圆面积}
用换元法求椭圆 $\frac{x^2}{4}+y^2=1$ 内部区域的面积。
\end{exercise}
\begin{exercise}{677}{基础·抛物线面积}
求 $y=2x-x^2$ 与 $x$ 轴围成区域的面积。
\end{exercise}
\begin{exercise}{678}{常规·分段面积}
求 $y=x^2$ 与 $y=\frac{1}{4}$ 在 $0\le x\le1$ 之间所夹区域的总面积。
\end{exercise}
\begin{exercise}{679}{基础·两抛物线面积}
求 $y=x^2$ 与 $y=2-x^2$ 围成区域的面积。
\end{exercise}
\begin{exercise}{680}{常规·面积随参数变化}
设 $a>0$，求 $y=ax$ 与 $y=x^2$ 围成区域的面积；面积为 $\frac{1}{6}$ 时求 $a$。
\end{exercise}
\begin{exercise}{681}{基础·极坐标面积}
用极坐标面积公式求曲线 $r=2\cos\theta$ 围成圆盘的面积，其中 $-\frac{\pi}{2}\le\theta\le\frac{\pi}{2}$。
\end{exercise}
\begin{exercise}{682}{常规·极坐标面积}
求心形线 $r=1+\cos\theta$（$0\le\theta\le2\pi$）围成区域的面积。
\end{exercise}
\begin{exercise}{683}{常规·参数曲线面积}
摆线一拱为 $x=t-\sin t$、$y=1-\cos t$，$0\le t\le2\pi$。求它与 $x$ 轴围成区域的面积。
\end{exercise}
\begin{exercise}{684}{基础·参数曲线面积}
曲线由 $x=t^2$、$y=t$（$0\le t\le1$）给出。求该曲线与 $x$ 轴、$x=1$ 围成区域的面积。
\end{exercise}
\begin{exercise}{685}{基础·圆盘法体积}
求 $y=\sqrt x$、$x=0$、$x=4$、$y=0$ 围成区域绕 $x$ 轴旋转一周的体积。
\end{exercise}
\begin{exercise}{686}{常规·圆环法体积}
求 $y=x$ 与 $y=x^2$ 围成区域绕 $x$ 轴旋转一周的体积，注明内、外半径。
\end{exercise}
\begin{exercise}{687}{基础·无界旋转体积}
求 $D=\{(x,y):x\ge0,0\le y\le e^{-x}\}$ 绕 $x$ 轴旋转所得体积，用反常积分表示。
\end{exercise}
\begin{exercise}{688}{基础·柱壳法体积}
求 $y=e^{-x^2}$、$x=-1$、$x=0$ 与 $x$ 轴围成区域绕 $y$ 轴旋转的体积，注意柱壳半径为 $-x$。
\end{exercise}
\begin{exercise}{689}{基础·圆盘旋转}
求圆盘 $x^2+y^2\le4$ 绕 $x$ 轴旋转所得球体的体积，先写出垂直于 $x$ 轴的截面面积。
\end{exercise}
\begin{exercise}{690}{基础·对数区域旋转}
求 $y=\ln x$、$y=0$ 与 $x=e$ 围成区域绕 $y$ 轴旋转的体积。
\end{exercise}
\begin{exercise}{691}{常规·反常积分体积}
求 $y=xe^{-x}$（$x\ge0$）与 $x$ 轴间区域绕 $x$ 轴旋转所得体积。
\end{exercise}
\begin{exercise}{692}{基础·半圆旋转体积}
求 $y=\sqrt{1-x^2}$（$-1\le x\le1$）与 $x$ 轴围成区域绕 $x$ 轴旋转的体积。
\end{exercise}
\begin{exercise}{693}{常规·平移旋转轴}
求矩形 $0\le x\le2$、$0\le y\le1$ 绕直线 $y=2$ 旋转一周的体积，注明内、外半径。
\end{exercise}
\begin{exercise}{694}{常规·切线区域旋转}
曲线 $y=x^2$ 在 $(1,1)$ 处的切线为 $y=2x-1$。求由该切线、曲线与 $y$ 轴围成区域绕 $y$ 轴旋转一周的体积，提示：使用柱壳法，$0\le x\le1$。
\end{exercise}
\begin{exercise}{695}{常规·累积体积}
设 $V(a)$ 为 $y=e^{-x}$、$y=0$、$x=0$、$x=a>0$ 围成区域绕 $x$ 轴旋转的体积。求 $V(a)$，并求满足 $V(a)=\frac{\pi}{4}$ 的 $a$。
\end{exercise}
\begin{exercise}{696}{常规·旋转曲面面积}
求直线段 $y=x$（$0\le x\le1$）绕 $x$ 轴旋转所得曲面的侧面积。
\end{exercise}
\begin{exercise}{697}{常规·旋转曲面面积}
求曲线 $y=\sqrt{2x+1}$（$0\le x\le1$）绕 $x$ 轴旋转所得曲面的面积。
\end{exercise}
\begin{exercise}{698}{基础·直角坐标弧长}
求曲线 $y=x^{\frac{3}{2}}$（$0\le x\le4$）的弧长。
\end{exercise}
\begin{exercise}{699}{基础·三角对数弧长}
求曲线 $y=\ln\cos x$（$0\le x\le\frac{\pi}{3}$）的弧长。
\end{exercise}
\begin{exercise}{700}{常规·变限函数弧长}
求曲线 $y=\int_0^x\tan t\dd t$（$0\le x\le\frac{\pi}{4}$）的弧长。
\end{exercise}
\begin{exercise}{701}{基础·参数曲线弧长}
求 $x=2\cos t$、$y=2\sin t$（$0\le t\le\frac{\pi}{2}$）所表示圆弧的长度。
\end{exercise}
\begin{exercise}{702}{基础·极坐标弧长}
求对数螺线 $r=e^\theta$（$0\le\theta\le\pi$）的弧长。
\end{exercise}
\begin{exercise}{703}{常规·参数弧长}
求曲线 $x=t^2$、$y=\frac23t^3$（$0\le t\le1$）的弧长。
\end{exercise}
\begin{exercise}{704}{常规·封闭区域周长}
求 $y=x$、$y=0$ 与 $x=1$ 围成三角形的周长，并用弧长积分核验斜边长度。
\end{exercise}
\begin{exercise}{705}{基础·曲率半径}
利用 $\rho=\frac{(1+(y')^2)^{\frac{3}{2}}}{|y''|}$，求抛物线 $y=x^2$ 在原点的曲率半径。
\end{exercise}
\begin{exercise}{706}{基础·轴外柱壳体积}
求矩形 $0\le x\le1$、$0\le y\le2$ 绕直线 $x=2$ 旋转一周的体积。
\end{exercise}
\begin{exercise}{707}{基础·绝对值区域旋转}
求 $y=1-|x|$ 与 $x$ 轴围成区域绕 $x$ 轴旋转所得体积。
\end{exercise}
\begin{exercise}{708}{常规·已知截面求体积}
某立体在 $0\le x\le2$ 内垂直于 $x$ 轴的截面均为正方形，边长为 $\sqrt{4-x^2}$。求立体体积。
\end{exercise}
\begin{exercise}{709}{基础·平移轴旋转变式}
求 $y=x$、$y=0$ 与 $x=1$ 围成区域绕直线 $y=2$ 旋转的体积。
\end{exercise}
\begin{exercise}{710}{基础·旋转体积变式}
求 $y=x^2$ 与 $y=2x$ 围成区域绕 $x$ 轴旋转的体积。
\end{exercise}
\begin{exercise}{711}{常规·有限旋转体积}
求 $D=\{(x,y):x\ge1,0\le y\le\frac{1}{x}\}$ 绕 $x$ 轴旋转所得体积，先截断到 $x=R$ 再取极限。
\end{exercise}
\subsection{定积分在物理学上的应用}
\begin{exercise}{712}{基础·速度与路程}
质点在 $0\le t\le2$ 内的速度为 $v(t)=3t^2+1$ 米/秒。求这段时间的路程。
\end{exercise}
\begin{exercise}{713}{常规·缆绳做功}
长 $10$ 米的均匀缆绳竖直悬挂，每米重 $20$ 牛。把缆绳全部缓慢拉到上端，求克服重力所做的功。
\end{exercise}
\begin{exercise}{714}{基础·矩形闸门压力}
竖直矩形闸门宽 $2$ 米、高 $3$ 米，上边位于水面。水密度为 $\rho$、重力加速度为 $g$。用水平薄条积分求单侧静水总压力。
\end{exercise}
\begin{exercise}{715}{基础·变力做功}
弹簧伸长 $x$ 米时弹力大小为 $100x$ 牛。将弹簧从自然长度拉长到 $0.2$ 米，求克服弹力所做的功。
\end{exercise}
\begin{exercise}{716}{基础·变密度细棒质量}
长 $2$ 米的细棒在位置 $x$ 米处的线密度为 $1+x$ 千克/米（$0\le x\le2$）。求细棒质量。
\end{exercise}
\begin{exercise}{717}{常规·圆柱容器抽水功}
半径 $1$ 米、高 $2$ 米的竖直圆柱水箱盛满水。求将水全部抽至箱顶所做的功，水密度为 $1000$ 千克/立方米，重力加速度为 $g$。
\end{exercise}
\begin{exercise}{718}{常规·均匀薄片形心}
均匀薄片为 $0\le x\le1$、$0\le y\le x$。求其面积及形心横坐标，提示：$\bar x=\frac{\int_0^1x\cdot x\dd x}{\int_0^1x\dd x}$。
\end{exercise}
\begin{exercise}{719}{基础·变力做功变式}
物体沿 $x$ 轴正向移动，所受力为 $F(x)=2x+3$ 牛。求物体从 $x=0$ 米移到 $x=2$ 米时该力所做的功。
\end{exercise}
\begin{exercise}{720}{常规·三角形水压小综合}
竖直等腰三角形薄板高 $h$、底边长 $b$，顶点在水面，底边水平且位于顶点下方。水密度为 $\rho$、重力加速度为 $g$。求薄板单侧静水总压力，提示：深度 $y$ 处的宽度为 $\frac{by}{h}$。
\end{exercise}
